  \section{Lineare Regression}
  \setcounter{aufgabe}{0}

  \aufgabe{
Gegeben seien folgende Testdatenpunkte aus $\mathbb{R} \times \mathbb{R}$:
\[
(0,\,4),\quad (1,\,5),\quad (2,\,7),\quad (3,\,12).
\]
Weiter sei die Regressionsfunktion $f: \mathbb{R}\to \mathbb{R}$ mit
\[
f(\mathbf{x}) = u_0 + \mathbf{b}\cdot\mathbf{x}, \qquad u_0 = 2,\ \mathbf{b} = (3)
\]
gegeben, also $f(x) = 2 + 3x$.

Bestimmen Sie die Summe der quadrierten Fehler $E$.
}

\loesung{
Es gilt mit $m$ Datenpunkten:
\[
E = \sum_{i=1}^{m}\bigl(y_i - f(\mathbf{x}_i)\bigr)^2.
\]

Berechnung der Vorhersagen und Fehler ($m=4$):
\begin{itemize}
  \item $\hat{y}=f(0) = 2$, \quad $y=4$, \quad $(y-\hat{y})^2 = 4$
  \item $\hat{y}=f(1) = 5$, \quad $y=5$, \quad $(y-\hat{y})^2 = 0$
  \item $\hat{y}=f(2) = 8$, \quad $y=7$, \quad $(y-\hat{y})^2 = 1$
  \item $\hat{y}=f(3) = 11$, \quad $y=12$, \quad $(y-\hat{y})^2 = 1$
\end{itemize}

\[
E = 4+0+1+1 = 6.
\]
}


\aufgabe{
Gegeben seien die Testdatenpunkte im Format $((x_1,x_2,x_3)^{\mathsf{T}},\,y)$ mit
Eingangsdimension $n=3$:
\[
((4,0,0)^{\mathsf{T}},\,8),\quad ((2,1,1)^{\mathsf{T}},\,5),\quad ((1,1,1)^{\mathsf{T}},\,4),\quad ((0,0,2)^{\mathsf{T}},\,3).
\]
Weiter sei die Regressionsfunktion $f: \mathbb{R}^3 \to \mathbb{R}$ mit
\[
f(\mathbf{x}) = u_0 + \mathbf{b}\cdot\mathbf{x}, \qquad
u_0 = 1,\ \mathbf{b} = (2,-1,1)^{\mathsf{T}}
\]
gegeben.

Bestimmen Sie die Summe der quadrierten Fehler $E$.
}

\loesung{
Es gilt mit $m$ Datenpunkten:
\[
E = \sum_{i=1}^{m}\bigl(y_i - f(\mathbf{x}_i)\bigr)^2.
\]

Berechnung der Vorhersagen und Fehler ($m=4$):
\begin{itemize}
  \item $\hat{y}=f(4,0,0) = 1+8-0+0 = 9$, \quad $y=8$, \quad $(y-\hat{y})^2 = 1$
  \item $\hat{y}=f(2,1,1) = 1+4-1+1 = 5$, \quad $y=5$, \quad $(y-\hat{y})^2 = 0$
  \item $\hat{y}=f(1,1,1) = 1+2-1+1 = 3$, \quad $y=4$, \quad $(y-\hat{y})^2 = 1$
  \item $\hat{y}=f(0,0,2) = 1+0-0+2 = 3$, \quad $y=3$, \quad $(y-\hat{y})^2 = 0$
\end{itemize}

\[
E = 1+0+1+0 = 2.
\]
}

\aufgabe{
Gegeben seien die folgenden Testdatenpaare $(\mathbf{x},y)$ mit
$\mathbf{x}\in\mathbb{R}^5$ und $y\in\mathbb{R}$:
\[
((1,0,2,1,0)^{\mathsf{T}},\,5),\quad
((2,1,0,0,1)^{\mathsf{T}},\,5),\quad
((0,2,1,1,1)^{\mathsf{T}},\,4),
\]
\[
((1,1,1,2,0)^{\mathsf{T}},\,7),\quad
((3,0,1,0,2)^{\mathsf{T}},\,0).
\]

Weiter sei die lineare Regressionsfunktion
$f:\mathbb{R}^5\to\mathbb{R}$ gegeben durch
\[
f(\mathbf{x}) = u_0 + \mathbf{b} \cdot \mathbf{x},
\]
wobei
\[
u_0 = 3 \qquad\text{und}\qquad \mathbf{b} =(1,2,-1,1,-2)^{\mathsf{T}}
\]

Bestimmen Sie die Summe der quadrierten Fehler.
}


\loesung{
Es gilt mit $m$ Datenpunkten:
\[
E = \sum_{i=1}^{m}\bigl(y_i - f(\mathbf{x}_i)\bigr)^2.
\]

Berechnung der Vorhersagen und Fehler ($m=5$):
\begin{itemize}
  \item $\hat{y}=f(1,0,2,1,0) = 3+1+0-2+1-0 = 3$, \quad $y=5$, \quad $(y-\hat{y})^2 = 4$
  \item $\hat{y}=f(2,1,0,0,1) = 3+2+2-0+0-2 = 5$, \quad $y=5$, \quad $(y-\hat{y})^2 = 0$
  \item $\hat{y}=f(0,2,1,1,1) = 3+0+4-1+1-2 = 5$, \quad $y=4$, \quad $(y-\hat{y})^2 = 1$
  \item $\hat{y}=f(1,1,1,2,0) = 3+1+2-1+2-0 = 7$, \quad $y=7$, \quad $(y-\hat{y})^2 = 0$
  \item $\hat{y}=f(3,0,1,0,2) = 3+3+0-1+0-4 = 1$, \quad $y=0$, \quad $(y-\hat{y})^2 = 1$
\end{itemize}

\[
E = 4+0+1+0+1 = 6.
\]
}


\aufgabe{
Gegeben seien die Datenpunkte
\[
(x_1,y_1)=(0,1),\qquad
(x_2,y_2)=(1,5),\qquad
(x_3,y_3)=(2,3).
\]

Gesucht ist die Regressionsgerade
\[
y = u_0 + u_1 x.
\]

\begin{enumerate}[label=\alph*)]
  \item Stellen Sie die Matrix $A$ und den Vektor $\mathbf{y}$ auf.
  \item Berechnen Sie $A^{\mathsf{T}} A$.
  \item Invertieren Sie $A^{\mathsf{T}} A$.
  \item Bestimmen Sie den Gewichtsvektor $\mathbf{u} = (A^{\mathsf{T}} A)^{-1} A^{\mathsf{T}} \mathbf{y}$.
\end{enumerate}
}

\loesung{
\begin{enumerate}[label=\alph*)]
  \item Jede Zeile enthält den Offset $1$ und den Merkmalswert $x$:
  \[
  A = \begin{pmatrix} 1 & 0 \\ 1 & 1 \\ 1 & 2 \end{pmatrix},
  \qquad
  \mathbf{y} = \begin{pmatrix} 1 \\ 5 \\ 3 \end{pmatrix}.
  \]

  \item 
  \[
  A^{\mathsf{T}} A =
  \begin{pmatrix} 1 & 1 & 1 \\ 0 & 1 & 2 \end{pmatrix}
  \begin{pmatrix} 1 & 0 \\ 1 & 1 \\ 1 & 2 \end{pmatrix}
  =
  \begin{pmatrix} 3 & 3 \\ 3 & 5 \end{pmatrix}.
  \]

  \item Determinante: $\det(A^{\mathsf{T}} A) = 3\cdot 5 - 3\cdot 3 = 6$.
  Für eine $2\times 2$-Matrix gilt
  $\begin{pmatrix} a & b \\ c & d \end{pmatrix}^{-1}
  = \frac{1}{ad-bc}\begin{pmatrix} d & -b \\ -c & a \end{pmatrix}$,
  also
  \[
  (A^{\mathsf{T}} A)^{-1}
  = \frac{1}{6}\begin{pmatrix} 5 & -3 \\ -3 & 3 \end{pmatrix}.
  \]

  \item 
  \[
  A^{\mathsf{T}} \mathbf{y}
  = \begin{pmatrix} 1 & 1 & 1 \\ 0 & 1 & 2 \end{pmatrix}
    \begin{pmatrix} 1 \\ 5 \\ 3 \end{pmatrix}
  = \begin{pmatrix} 1+5+3 \\ 0+5+6 \end{pmatrix}
  = \begin{pmatrix} 9 \\ 11 \end{pmatrix}.
  \]
  \[
  \mathbf{u}
  = \frac{1}{6}\begin{pmatrix} 5 & -3 \\ -3 & 3 \end{pmatrix}
    \begin{pmatrix} 9 \\ 11 \end{pmatrix}
  = \frac{1}{6}\begin{pmatrix} 45-33 \\ -27+33 \end{pmatrix}
  = \frac{1}{6}\begin{pmatrix} 12 \\ 6 \end{pmatrix}
  = \begin{pmatrix} 2 \\ 1 \end{pmatrix}.
  \]
  Das Modell lautet also $f(x) = 2 + x$.
\end{enumerate}
}


\aufgabe{
Gegeben seien folgende Beobachtungen aus $\mathbb{R}^2\times \mathbb{R}$ im Format $(x_1, x_2, y)$
\[
 \bigl\{(1,1,3),\,(2,1,6),\,(1,2,7),\,(2,2,6)\bigr\}.
\]
Gesucht ist die Regressionsebene 
\[f:\mathbb{R}^2 \to \mathbb{R}, \mathbf{x} \mapsto u_0+u_1x_1 + u_2 x_2 \] 
mit Gewichtsvektor $\mathbf{u} = (u_0, u_1, u_2)^{\mathsf{T}} \in \mathbb{R}^3$.\par\medskip

\begin{enumerate}[label=\alph*)]
  \item Stellen Sie die Matrix $A \in \mathbb{R}^{4\times 3}$ und den Zielvektor $\mathbf{y} \in \mathbb{R}^4$ auf, sodass gilt $\mathbf{y} = A\mathbf{u}$.
  \item Berechnen Sie $A^{\mathsf{T}} A$.
  \item Bestimmen Sie die Inverse $\bigl(A^{\mathsf{T}} A\bigr)^{-1}$.
  \item Berechnen Sie den Gewichtsvektor $\mathbf{u} = \bigl(A^{\mathsf{T}} A\bigr)^{-1} A^{\mathsf{T}} \mathbf{y}$.
\end{enumerate}
}

\loesung{
\begin{enumerate}[label=\alph*)]
  \item Jede Zeile enthält den Offset $1$ sowie die Merkmale $x_1, x_2$:
  \[
  A = \begin{pmatrix}
  1 & 1 & 1 \\
  1 & 2 & 1 \\
  1 & 1 & 2 \\
  1 & 2 & 2
  \end{pmatrix},
  \qquad
  \mathbf{y} = \begin{pmatrix} 3 \\ 6 \\ 7 \\ 6 \end{pmatrix}.
  \]

  \item Die Einträge ergeben sich aus den Spaltensummen und Kreuzsummen
  ($n=4$, $\sum x_1 = 6$, $\sum x_2 = 6$, $\sum x_1^2 = 10$, $\sum x_2^2 = 10$, $\sum x_1 x_2 = 9$):
  \[
  A^{\mathsf{T}} A =
  \begin{pmatrix}
  4 & 6 & 6 \\
  6 & 10 & 9 \\
  6 & 9 & 10
  \end{pmatrix}.
  \]

  \item Entwicklung nach der ersten Zeile:
  \[
  \det(A^{\mathsf{T}} A)
  = 4(10\cdot 10 - 9\cdot 9) - 6(6\cdot 10 - 9\cdot 6) + 6(6\cdot 9 - 10\cdot 6)
  = 4\cdot 19 - 6\cdot 6 + 6\cdot(-6) = 4.
  \]
  Mit der Adjunkten (Kofaktormatrix) folgt:
  \[
  (A^{\mathsf{T}} A)^{-1}
  = \frac{1}{4}
  \begin{pmatrix}
  19 & -6 & -6 \\
  -6 & 4 & 0 \\
  -6 & 0 & 4
  \end{pmatrix}.
  \]

  \item Zunächst
  \[
  A^{\mathsf{T}} \mathbf{y} =
  \begin{pmatrix}
  3+6+7+6 \\
  3+12+7+12 \\
  3+6+14+12
  \end{pmatrix}
  = \begin{pmatrix} 22 \\ 34 \\ 35 \end{pmatrix}.
  \]
  Damit
  \[
  \mathbf{u}
  = \frac{1}{4}
  \begin{pmatrix}
  19 & -6 & -6 \\
  -6 & 4 & 0 \\
  -6 & 0 & 4
  \end{pmatrix}
  \begin{pmatrix} 22 \\ 34 \\ 35 \end{pmatrix}
  = \frac{1}{4}\begin{pmatrix} 418-204-210 \\ -132+136 \\ -132+140 \end{pmatrix}
  = \frac{1}{4}\begin{pmatrix} 4 \\ 4 \\ 8 \end{pmatrix}
  = \begin{pmatrix} 1 \\ 1 \\ 2 \end{pmatrix}.
  \]
  Das Modell lautet also $f(x_1,x_2) = 1 + x_1 + 2x_2$.
\end{enumerate}
}





